\documentclass[11pt,a4paper]{article}
\usepackage[utf8]{inputenc}
\usepackage[margin=1in]{geometry}
\usepackage{amsmath,amssymb,amsthm}
\usepackage{listings}
\usepackage{xcolor}
\usepackage{hyperref}

\newtheorem{theorem}{Theorem}
\newtheorem{lemma}[theorem]{Lemma}

\lstset{
  language=C++,
  basicstyle=\ttfamily\small,
  keywordstyle=\color{blue}\bfseries,
  commentstyle=\color{gray},
  stringstyle=\color{red},
  numbers=left,
  numberstyle=\tiny\color{gray},
  breaklines=true,
  frame=single,
  tabsize=4,
  showstringspaces=false
}

\title{IOI 2014 -- Gondola}
\author{}
\date{}

\begin{document}
\maketitle

\section{Problem Summary}

A gondola system has $n$ gondolas arranged in a circle, originally numbered $1, 2, \ldots, n$. Over time, gondolas may break and be replaced: the $k$-th replacement gondola is numbered $n + k$. The circular order is preserved. There are three subtasks:
\begin{enumerate}
  \item \textbf{Valid}: Determine if a given sequence is a valid gondola configuration.
  \item \textbf{Replacement}: Given a valid sequence, find the sequence of broken gondolas.
  \item \textbf{Count}: Count the number of valid replacement sequences yielding a given final configuration (modulo $10^9 + 9$).
\end{enumerate}

\section{Solution}

\subsection{Subtask 1: Validity}

A sequence $s_0, s_1, \ldots, s_{n-1}$ is valid if and only if:
\begin{enumerate}
  \item No value appears more than once.
  \item All \emph{original} gondolas ($s_i \le n$) are consistent with a single rotation. If $s_i \le n$, the rotation offset is $(s_i - 1 - i) \bmod n$. All such values must agree on the same offset.
\end{enumerate}

\subsection{Subtask 2: Replacement Sequence}

Determine the rotation offset from any original gondola present (default offset $0$ if none exists). Then, for each position with a replacement gondola ($s_i > n$), record $(\text{value}, \text{position})$. Sort these pairs by value. Process them in order: the first replacement at a position replaces the original gondola there; subsequent replacements at that position replace the previous replacement.

\subsection{Subtask 3: Counting}

Let $\text{maxVal} = \max(s_i)$. There are $\text{maxVal} - n$ total replacements. Sort the replacement positions by their final values. Let these sorted values be $v_1 < v_2 < \cdots < v_r$. Between consecutive ``pinned'' replacements, each intermediate replacement can go to any of the currently-unresolved positions. Specifically, if $c$ positions are still unresolved, and there is a gap of $g$ intermediate replacements, we multiply the answer by $c^g$. If no original gondola is present, any of $n$ rotations is valid, contributing a factor of $n$.

\section{C++ Implementation}

\begin{lstlisting}
#include <bits/stdc++.h>
using namespace std;

const int MOD = 1e9 + 9;

long long power(long long base, long long exp, long long mod) {
    long long result = 1;
    base %= mod;
    while (exp > 0) {
        if (exp & 1) result = result * base % mod;
        base = base * base % mod;
        exp >>= 1;
    }
    return result;
}

int valid(int n, int inputSeq[]) {
    set<int> seen;
    int offset = -1;
    for (int i = 0; i < n; i++) {
        if (seen.count(inputSeq[i])) return 0;
        seen.insert(inputSeq[i]);
        if (inputSeq[i] <= n) {
            int off = ((inputSeq[i] - 1) - i % n + n) % n;
            if (offset == -1) offset = off;
            else if (offset != off) return 0;
        }
    }
    return 1;
}

int replacement(int n, int inputSeq[], int replacementSeq[]) {
    int offset = 0;
    for (int i = 0; i < n; i++) {
        if (inputSeq[i] <= n) {
            offset = ((inputSeq[i] - 1) - i + n) % n;
            break;
        }
    }
    vector<pair<int,int>> reps;
    for (int i = 0; i < n; i++)
        if (inputSeq[i] > n)
            reps.push_back({inputSeq[i], i});
    sort(reps.begin(), reps.end());

    int len = 0, nextId = n + 1;
    for (auto &[val, pos] : reps) {
        int orig = (pos + offset) % n + 1;
        replacementSeq[len++] = orig;
        for (int id = nextId; id < val; id++)
            replacementSeq[len++] = id;
        nextId = val + 1;
    }
    return len;
}

int countReplacement(int n, int inputSeq[]) {
    if (!valid(n, inputSeq)) return 0;

    int offset = -1;
    for (int i = 0; i < n; i++) {
        if (inputSeq[i] <= n) {
            offset = ((inputSeq[i] - 1) - i + n) % n;
            break;
        }
    }

    vector<pair<int,int>> reps;
    for (int i = 0; i < n; i++)
        if (inputSeq[i] > n)
            reps.push_back({inputSeq[i], i});
    sort(reps.begin(), reps.end());

    int cntUnknown = (int)reps.size();
    long long ans = 1;
    if (offset == -1 && cntUnknown > 0)
        ans = n;  // any rotation is valid

    long long prev = n;
    for (int i = 0; i < (int)reps.size(); i++) {
        long long cur = reps[i].first;
        long long gap = cur - prev - 1;
        ans = ans % MOD * power(cntUnknown, gap, MOD) % MOD;
        cntUnknown--;
        prev = cur;
    }
    return (int)(ans % MOD);
}
\end{lstlisting}

\section{Complexity Analysis}

\begin{itemize}
  \item \textbf{Subtask 1}: $O(n \log n)$ (set operations).
  \item \textbf{Subtask 2}: $O(n \log n + (\max s_i - n))$.
  \item \textbf{Subtask 3}: $O(n \log n + n \log(\max s_i))$ (sorting + modular exponentiation).
  \item \textbf{Space}: $O(n)$.
\end{itemize}

\end{document}
